\documentclass[11pt,reqno]{amsart}

\usepackage[utf8]{inputenc}
\usepackage{amsmath,amssymb,mathtools}
\usepackage{microtype}
\usepackage{xcolor}
\usepackage{hyperref}

\hypersetup{
  colorlinks=true,
  linkcolor=blue!55!black,
  citecolor=blue!55!black,
  urlcolor=blue!65!black
}

\newtheorem{theorem}{Theorem}[section]
\newtheorem{proposition}[theorem]{Proposition}
\newtheorem{lemma}[theorem]{Lemma}
\newtheorem{corollary}[theorem]{Corollary}
\theoremstyle{remark}
\newtheorem{remark}[theorem]{Remark}

\DeclareMathOperator{\erf}{erf}
\DeclareMathOperator{\erfc}{erfc}

\newcommand{\R}{\mathbb{R}}
\newcommand{\Hads}{\mathcal{H}_{\mathrm{ADS}}}
\newcommand{\Haus}{\mathcal{H}}
\newcommand{\Pgamma}{P_{\gamma}}

\title[An explicit counterexample to a Huisken-energy bound]
{An explicit counterexample to a conjectured dimension-free
Huisken-energy bound}

% Replace the following line with the author's name before submission.
\author{ChatGPT 5.6 Sol}
% \address{Department, University, Postal address}
% \email{author@example.edu}

\subjclass[2020]{53E10, 52A20, 60D05}
\keywords{mean curvature flow, Huisken energy, Gaussian perimeter,
convex body, counterexample}

\begin{document}

\begin{abstract}
\textbf{Relation to the problem list.}
This manuscript addresses Question~5 (L5), ``Gaussian boundary area of
convex bodies is below two.''  It gives a full counterexample to the literal
dimension-free conjecture, including a smooth strictly convex example.  The
disproof itself is not new: after the correct normalization it follows from
Ball's 1993 cube estimate; the contribution here is an explicit
low-dimensional exact certificate.

Angenent, Daskalopoulos, and \v{S}e\v{s}um conjectured that the
Gaussian-weighted boundary area
\[
 (4\pi)^{-(d-1)/2}\int_{\partial C}e^{-|x|^2/4}\,
 d\Haus^{d-1}(x)
\]
is strictly smaller than \(2\) for every closed convex set
\(C\subset\R^d\) with nonempty interior.  We record a completely
explicit counterexample in ambient dimension \(9\):
\[
 Q=\left[-\frac{12}{5},\frac{12}{5}\right]^9.
\]
An exact calculation, certified using explicit rational bounds, gives
\[
 \Hads(Q)>\frac{1002087}{500000}>2.
\]
We also show that the same violation
persists for smooth convex bodies whose principal curvatures are everywhere
positive.  After a dilation, the functional is
\(\sqrt{2\pi}\) times the standard Gaussian perimeter.  Accounting for this
normalization, the failure of a dimension-free bound follows directly from
Ball's 1993 cube estimate.  The purpose of this note is to spell out its
consequence for the stated conjecture and to provide a low-dimensional exact
certificate.
\end{abstract}

\maketitle

\section{Introduction}

Let \(d\ge2\).  For a closed convex set \(C\subset\R^d\) with nonempty
interior, set
\begin{equation}\label{eq:ADS-functional}
 \Hads(C)
 :=(4\pi)^{-(d-1)/2}
   \int_{\partial C}e^{-|x|^2/4}\,d\Haus^{d-1}(x).
\end{equation}
In the notation of Angenent--Daskalopoulos--\v{S}e\v{s}um
\cite[Conjecture~I]{ADS}, the ambient dimension is \(d=n+1\) and
\(\Hads\) is their Huisken energy.  Their conjecture asserts
\begin{equation}\label{eq:conjecture}
 \Hads(C)<2
 \qquad\text{for every closed convex \(C\subset\R^d\) with nonempty
 interior.}
\end{equation}
The phase space in \cite[\S1.2]{ADS} imposes no regularity or boundedness
condition beyond convexity, so polytopes are included.

Our main result is the following.

\begin{theorem}\label{thm:main}
In ambient dimension \(d=9\), the compact centrally symmetric convex body
\[
 Q=\left[-\frac{12}{5},\frac{12}{5}\right]^9
\]
satisfies
\begin{equation}\label{eq:main-bound}
 \Hads(Q)>\frac{1002087}{500000}>2.
\end{equation}
Consequently, Conjecture~\eqref{eq:conjecture} is false already for
hypersurface dimension \(n=8\).
\end{theorem}

The corners of \(Q\) are not responsible for the failure.

\begin{theorem}\label{thm:smooth}
There exists a compact convex body \(C\subset\R^9\) with smooth boundary
and everywhere positive principal curvatures such that
\(\Hads(C)>2\).
\end{theorem}

\section{The weighted boundary area of a cube}

For \(a>0\) and an integer \(d\ge2\), write
\[
 Q_d(a):=[-a,a]^d.
\]
Recall that
\[
 \erf(s):=\frac{2}{\sqrt\pi}\int_0^s e^{-t^2}\,dt,
 \qquad
 \erfc(s):=1-\erf(s).
\]

\begin{proposition}[Cube formula]\label{prop:cube-formula}
For every \(a>0\) and \(d\ge2\),
\begin{equation}\label{eq:cube-formula}
 \Hads\bigl(Q_d(a)\bigr)
 =2d\,e^{-a^2/4}\erf(a/2)^{d-1}.
\end{equation}
\end{proposition}

\begin{proof}
The intersections of distinct facets have zero \((d-1)\)-dimensional
Hausdorff measure.  On any one of the \(2d\) facets, one coordinate is
equal to \(a\) or \(-a\), and Fubini's theorem gives
\[
 \int_{[-a,a]^{d-1}}e^{-(a^2+|x'|^2)/4}\,dx'
 =e^{-a^2/4}
   \left(\int_{-a}^{a}e^{-t^2/4}\,dt\right)^{d-1}.
\]
The substitution \(t=2s\) yields
\[
 \int_{-a}^{a}e^{-t^2/4}\,dt
 =2\sqrt\pi\,\erf(a/2).
\]
Multiplying by \(2d(4\pi)^{-(d-1)/2}\) proves
\eqref{eq:cube-formula}.
\end{proof}

For later perspective, the same formula gives a self-contained version of
the classical cube-growth observation.

\begin{proposition}[No dimension-free upper bound]\label{prop:unbounded}
The values of \(\Hads\) on compact centrally symmetric cubes are unbounded
as the ambient dimension tends to infinity.
\end{proposition}

\begin{proof}
For an integer \(k\ge1\), let
\[
 d_k:=\left\lceil k e^{k^2}\right\rceil,
 \qquad Q_k:=[-2k,2k]^{d_k},
 \qquad q_k:=\erfc(k).
\]
The elementary Mills bound
\begin{equation}\label{eq:mills}
 \erfc(s)
 =\frac{2}{\sqrt\pi}\int_s^\infty e^{-t^2}\,dt
 <\frac{e^{-s^2}}{s\sqrt\pi}
 \qquad(s>0)
\end{equation}
follows by replacing \(1\) with \(t/s\) under the integral.  Since
\(d_k\le k e^{k^2}+1\), the elementary inequalities \(e>2\) and
\(\sqrt\pi>3/2\) imply
\[
 d_kq_k
 <\frac1{\sqrt\pi}+\frac{e^{-k^2}}{k\sqrt\pi}
 <\frac23+\frac13=1,
 \qquad
 q_k<\frac{e^{-1}}{\sqrt\pi}<\frac13<\frac12.
\]
For \(0\le q\le1/2\), one has
\(\log(1-q)\ge-2q\).  Therefore
\[
 (1-q_k)^{d_k-1}
 =\exp\bigl((d_k-1)\log(1-q_k)\bigr)>e^{-2}.
\]
Proposition~\ref{prop:cube-formula} and
\(d_ke^{-k^2}\ge k\) now give
\[
 \Hads(Q_k)
 =2d_ke^{-k^2}(1-q_k)^{d_k-1}
 >2e^{-2}k,
\]
which tends to infinity.
\end{proof}

\section{An exact certificate in dimension nine}

The proof of Theorem~\ref{thm:main} will use only the following rational
bounds.  Their somewhat explicit form is included so that the strict
inequality does not rest on floating-point computation.

\begin{lemma}[Rational certificate]\label{lem:certificate}
The following inequalities hold:
\begin{equation}\label{eq:certificate}
 e^{-36/25}>\frac{2369}{10000},
 \qquad
 \erf(6/5)>\frac{91}{100},
 \qquad
 \left(\frac{91}{100}\right)^8>\frac{47}{100}.
\end{equation}
\end{lemma}

\begin{proof}
We prove the three estimates in turn.

\smallskip
\noindent
\emph{The exponential estimate.}
Put \(y=36/25\).  For the terms \(t_j=y^j/j!\), one has
\(t_{j+1}/t_j\le y/7=36/175\) whenever \(j\ge6\).  Hence
\begin{align*}
 e^y
 &<\sum_{j=0}^{5}\frac{y^j}{j!}
   +\frac{1}{1-36/175}\frac{y^6}{6!}\\
 &=\frac{28647109619}{6787109375}
 <\frac{4221}{1000}.
\end{align*}
For completeness, the final difference is
\[
 \frac{4221}{1000}
 -\frac{28647109619}{6787109375}
 =\frac{10232423}{54296875000}>0.
\]
Taking reciprocals and clearing denominators once more gives
\[
 e^{-36/25}>\frac{1000}{4221}>\frac{2369}{10000},
 \qquad
 10000000-2369\cdot4221=451>0.
\]

\smallskip
\noindent
\emph{The error-function estimate.}
We use the elementary identity
\[
 \frac{22}{7}-\pi
 =\int_0^1\frac{x^4(1-x)^4}{1+x^2}\,dx>0,
\]
which follows by polynomial division and
\(\int_0^1(1+x^2)^{-1}\,dx=\pi/4\).  Thus \(\pi<22/7\), and hence
\begin{equation}\label{eq:pi-bound}
 \frac2{\sqrt\pi}>\frac{11281}{10000},
\end{equation}
because
\[
 11\cdot11281^2=1399870571
 <1400000000=14\cdot10000^2.
\]
The power series
\begin{equation}\label{eq:erf-series}
 \erf(x)=\frac2{\sqrt\pi}
 \sum_{j=0}^{\infty}
 \frac{(-1)^j x^{2j+1}}{j!(2j+1)}
\end{equation}
is alternating.  At \(x=6/5\), its term magnitudes decrease strictly,
since
\[
 \frac{(6/5)^{2j+3}}{(j+1)!(2j+3)}
 \bigg/
 \frac{(6/5)^{2j+1}}{j!(2j+1)}
 =\frac{36(2j+1)}{25(j+1)(2j+3)}<1;
\]
after clearing denominators, the denominator minus the numerator is
\(50j^2+53j+39>0\).  The partial sum ending with the negative term
\(j=7\) is therefore a strict lower bound.  Direct rational summation gives
\begin{align*}
 \sum_{j=0}^{7}
 \frac{(-1)^j(6/5)^{2j+1}}{j!(2j+1)}
 &=\frac{616091309997726}{763702392578125}\\
 &>\frac{8067}{10000}.
\end{align*}
The last comparison is certified by the positive cross-product difference
\[
 10000\cdot616091309997726
 -8067\cdot763702392578125
 =125899049525625.
\]
Combining this with \eqref{eq:pi-bound} and \eqref{eq:erf-series},
\[
 \erf(6/5)
 >\frac{11281\cdot8067}{10^8}
 =\frac{91003827}{10^8}
 >\frac{91}{100}.
\]

\smallskip
\noindent
\emph{The power estimate.}
Finally,
\[
 \left(\frac{91}{100}\right)^4
 =\frac{68574961}{10^8}>\frac{6857}{10000},
\]
and therefore
\[
 \left(\frac{91}{100}\right)^8
 >\left(\frac{6857}{10000}\right)^2
 =\frac{47018449}{10^8}>\frac{47}{100}.
\]
This proves \eqref{eq:certificate}.
\end{proof}

\begin{proof}[Proof of Theorem~\ref{thm:main}]
Apply Proposition~\ref{prop:cube-formula} with \(d=9\) and \(a=12/5\).
Since \(a^2/4=36/25\) and \(a/2=6/5\), Lemma~\ref{lem:certificate}
gives
\begin{align*}
 \Hads(Q)
 &=18e^{-36/25}\erf(6/5)^8\\
 &>18\left(\frac{2369}{10000}\right)
       \left(\frac{47}{100}\right)
 =\frac{1002087}{500000}>2.
\end{align*}
The last inequality has the exact margin
\(2087/500000>0\).
\end{proof}

\section{Smooth strictly convex counterexamples}

We now prove that the inequality is not rescued by imposing smoothness or
strict convexity.  For every integer \(m\ge2\), define
\begin{equation}\label{eq:smooth-approximation}
 C_m:=\left\{x\in\R^9:
 \sum_{i=1}^9\left(\frac{5x_i}{12}\right)^{2m}
 +\frac1m\sum_{i=1}^9\left(\frac{5x_i}{12}\right)^2
 \le1\right\}.
\end{equation}

\begin{proposition}\label{prop:approximation}
Each \(C_m\) is a compact convex body with smooth boundary and everywhere
positive principal curvatures.  Moreover, \(C_m\) converges in Hausdorff
distance to \(Q=[-12/5,12/5]^9\).
\end{proposition}

\begin{proof}
Write \(z_i=5x_i/12\) and
\[
 F_m(z):=\sum_{i=1}^9z_i^{2m}+\frac1m\sum_{i=1}^9z_i^2.
\]
The change \(x=(12/5)z\) is a homothety, so smoothness, strict convexity,
and positivity of the principal curvatures may be checked in the
\(z\)-coordinates.
The Hessian of \(F_m\) in the \(z\)-coordinates is diagonal, and
\[
 \frac{\partial^2F_m}{\partial z_i^2}
 =2m(2m-1)z_i^{2m-2}+\frac2m>0.
\]
Thus \(F_m\) is strictly convex.  It is coercive, and \(F_m(0)=0\), so
its unit sublevel set is a compact convex body.  On the level set \(F_m=1\),
\[
 z\cdot\nabla F_m(z)
 =2m\sum_{i=1}^9z_i^{2m}+\frac2m\sum_{i=1}^9z_i^2>0,
\]
so the gradient never vanishes there.  Consequently
\(\partial C_m\) is smooth.  With the convention that the outward second
fundamental form of a convex hypersurface is positive, the second fundamental
form of the level set in \(z\)-space is the restriction of
\(D^2F_m/|\nabla F_m|\) to tangent vectors.  It is positive definite, and
the homothety back to \(x\)-space preserves that sign.  Therefore all
principal curvatures of \(\partial C_m\) are positive.

If \(x\in C_m\), then \(|5x_i/12|\le1\) for every \(i\), so \(C_m\subset Q\).
Conversely, fix \(0<\eta<1\).  For \(x\in(1-\eta)Q\),
\[
 F_m(5x/12)
 \le9(1-\eta)^{2m}+\frac9m(1-\eta)^2.
\]
The right-hand side is smaller than \(1\) for all sufficiently large \(m\).
Thus \((1-\eta)Q\subset C_m\subset Q\) eventually.  Every point of \(Q\)
lies within \(\eta\sup_{x\in Q}|x|\) of \((1-\eta)Q\subset C_m\), while
\(C_m\subset Q\); hence the Hausdorff distance tends to zero.
\end{proof}

\begin{proof}[Proof of Theorem~\ref{thm:smooth}]
Angenent--Daskalopoulos--\v{S}e\v{s}um prove that \(\Hads\) is continuous
under local Hausdorff convergence of closed convex sets
\cite[Lemma~3.6]{ADS}.  For the compact sets in
Proposition~\ref{prop:approximation}, ordinary Hausdorff convergence implies
local Hausdorff convergence.  Hence
\[
 \lim_{m\to\infty}\Hads(C_m)=\Hads(Q)>2.
\]
It follows that \(\Hads(C_m)>2\) for every sufficiently large \(m\).  Any
one such \(C_m\) has all the asserted properties.
\end{proof}

\begin{remark}
The approximation argument is deliberately qualitative: it proves the
existence of a smooth positively curved counterexample without claiming a
numerically optimal value of \(m\).  The polytope \(Q\) itself already lies
in the precise phase space of \cite{ADS} and is the exact counterexample to
the conjecture as stated.
\end{remark}

\begin{thebibliography}{99}

\bibitem{ADS}
S.~B. Angenent, P. Daskalopoulos, and N. \v{S}e\v{s}um,
\emph{Dynamics of convex mean curvature flow},
J. Reine Angew. Math. \textbf{825} (2025), 185--235,
\href{https://doi.org/10.1515/crelle-2025-0029}
{doi:10.1515/crelle-2025-0029};
preprint, \href{https://arxiv.org/abs/2305.17272}{arXiv:2305.17272}.

\bibitem{Ball}
K. Ball,
\emph{The reverse isoperimetric problem for Gaussian measure},
Discrete Comput. Geom. \textbf{10} (1993), 411--420,
\href{https://doi.org/10.1007/BF02573986}
{doi:10.1007/BF02573986}.

\bibitem{Nazarov}
F.~L. Nazarov,
\emph{On the maximal perimeter of a convex set in \(\R^n\) with respect to
a Gaussian measure},
in \emph{Geometric Aspects of Functional Analysis},
Lecture Notes in Math. 1807, Springer, Berlin, 2003, 169--187,
\href{https://doi.org/10.1007/978-3-540-36428-3_15}
{doi:10.1007/978-3-540-36428-3\_15}.

\end{thebibliography}

\end{document}
